% Chapter 1: A Century of Experiments
% 14 subsections, each a key experiment or result

Reinforcement learning emerged from a series of collisions between fields that did not know they were working on the same problem. Psychologists studying cats in puzzle boxes discovered that consequences shape behavior. Economists analyzing lotteries realized that the value of an outcome is not the outcome itself. Operations researchers solving logistics problems found that optimal decisions decompose recursively. And computer scientists building game-playing programs learned that a machine could improve its own evaluation function by playing against itself. These traditions developed largely in parallel for decades before converging, in the late twentieth century, into a unified computational framework for sequential decision-making under uncertainty. This chapter traces that convergence through fifteen experiments and results, each of which exposed a limitation in what came before and pointed toward what came next.

The sequence is roughly chronological, but the organizing principle is conceptual rather than calendrical. The chapter covers fifteen steps, spanning a century from Thorndike (1898) to Christiano (2017).

To aid comparison across these different traditions, this chapter adopts a single, modern set of notation for the core concepts of reinforcement learning, following current conventions \citep{sutton2018}. We denote a state by $s$, an action by $a$, a policy by $\pi$, the state-value function by $V(s)$, the action-value function by $Q(s,a)$, the reward by $r$, and the discount factor by $\gamma$. When presenting historical formalisms, we will often use the author's original notation first, and then explicitly map it to this modern framework. This sometimes requires a slight abuse of notation, as the original concepts may not be perfect analogues, but it serves the goal of creating a unified conceptual narrative.

\subsection{Thorndike's Cats (1898)}

Comparative psychology in the 1890s sought to measure animal intelligence through controlled experiment rather than anecdote. \citet{Thorndike1898} placed hungry cats inside latched wooden boxes with food visible outside and recorded escape latency across repeated trials. The experimental apparatus was simple: a wooden box with a door held shut by a latch mechanism that the cat could operate by pressing a lever, pulling a loop of string, or stepping on a treadle. The dependent variable was the time from placement inside the box to successful escape.

Trial 1 required 160 seconds of undirected clawing, biting, and squeezing through slats. By trial 24 the cat pressed the latch directly in 7 seconds. The learning curve was smooth rather than stepwise, ruling out sudden insight and suggesting incremental strengthening of stimulus-response associations.\footnote{Thorndike's escape-time data followed an approximate power law $T(n) \propto n^{-\beta}$ rather than an exponential decay, a pattern later termed the ``power law of practice.'' Whether animal and human learning curves are fundamentally power-law or a mixture of exponentials remains debated in the quantitative learning literature.} Thorndike codified this pattern as the Law of Effect:

\begin{quote}\itshape
Of several responses made to the same situation, those which are accompanied or closely followed by satisfaction to the animal will, other things being equal, be more firmly connected with the situation, so that, when it recurs, they will be more likely to recur \citep{Thorndike1898}.
\end{quote}

In modern terms, the box configuration is the state $s$, the motor action (pressing, pulling, clawing) is $a$, and escape followed by food is the reward $r$.\footnote{This is a slight abuse of notation, as Thorndike's ``situation'' was not a formally defined state and his ``responses'' were not drawn from a well-defined action set. The mapping is conceptual, highlighting the analogy between the Law of Effect and the principles of modern reinforcement learning.} The ``stamping in'' of successful responses increases the probability of selecting that action given the state, a directional update rule that Q-learning would formalize nine decades later. The puzzle box was an ideal preparation for demonstrating the Law of Effect because it provided a single well-defined state, a discrete set of motor actions, and an unambiguous binary reward: escape followed by food, or continued confinement. Each trial reset the cat to the same initial condition, producing independent episodes with a clean reinforcement signal. A richer environment with delayed, partial, or ambiguous feedback would have obscured the incremental strengthening that Thorndike sought to measure. What Thorndike lacked was any formal apparatus for representing the learning process mathematically. That apparatus would come from a different tradition entirely.

\subsection{The Rescorla-Wagner Model (1972)}

Classical conditioning theory through the mid-twentieth century held that temporal pairing of stimuli was sufficient for associative learning. The blocking experiment overturned this view by demonstrating that prediction error, not co-occurrence, drives associative change.

\begin{quote}\itshape
Organisms only learn when events violate their expectations. Certain expectations are built up about the events following a stimulus complex; expectations initiated by that complex and its component stimuli are then only modified when consequent events disagree with the composite expectation \citep{RescorlaWagner1972}.
\end{quote}

The protocol proceeds in two phases. In Phase I, a rat learns that a tone predicts shock and develops a conditioned fear response to the tone alone. In Phase II, the compound stimulus of tone plus light is paired with the same shock. When the light is subsequently presented alone, no fear response occurs. The light is ``blocked'' because the tone already predicts the shock perfectly, leaving no prediction error to drive learning about the light. \citet{RescorlaWagner1972} formalized this observation with the following update rule:
\begin{equation}
\Delta V_i = \alpha_i \beta_j (\lambda_j - V_{\text{tot}})
\end{equation}
where $V_i$ is the associative strength of stimulus $i$, $\alpha_i$ is the salience of stimulus $i$, $\beta_j$ is the learning rate associated with unconditioned stimulus $j$, $\lambda_j$ is the maximum conditioning that stimulus $j$ supports, and $V_{\text{tot}} = \sum_k V_k$ is the total prediction from all stimuli present on the trial.\footnote{The Rescorla-Wagner update $\Delta V_i = \alpha_i \beta_j (\lambda_j - V_{\text{tot}})$ is mathematically identical to the Widrow-Hoff (1960) least mean squares (LMS) rule used in adaptive signal processing. Both minimize mean squared prediction error by stochastic gradient descent. The parallel development in psychology and engineering was independent.} Learning halts when $V_{\text{tot}} = \lambda_j$, because the prediction error $(\lambda_j - V_{\text{tot}})$ reaches zero. In the blocking design, $V_{\text{tone}}$ already equals $\lambda$ after Phase I, so $\Delta V_{\text{light}} = 0$ throughout Phase II. The blocking design isolates the prediction error mechanism with surgical precision: Phase I drives $V_{\text{tone}}$ to $\lambda$, so Phase II produces exactly zero error on every trial, and the model predicts exactly zero learning about the light. A design with partial or probabilistic reinforcement would still be consistent with the model but would not produce the sharp qualitative prediction (complete absence of conditioning to the light) that made the experimental test so compelling.

The Rescorla-Wagner model can be mapped directly to the temporal difference framework. The term $V_{\text{tot}}$ is the value of the current state, $V(s)$. The term $\lambda_j$ represents the maximum possible reward, which in a single-step episode is the reward itself, $r$. The associative strength of a single stimulus, $V_i$, corresponds to a component of the value function. With a discount factor $\gamma=0$, the update is driven by the prediction error $(\lambda_j - V_{\text{tot}})$, which is equivalent to the TD error, $\delta = r - V(s)$. The model is thus a special case of TD learning for single-step, terminal-reward tasks. The deeper point is that learning is driven by the discrepancy between expected and received outcomes rather than by the outcomes themselves, a principle that would become the conceptual foundation of all temporal difference methods.

\subsection{The St.\ Petersburg Paradox and the Reward Function (1738)}

The theory of decision under risk began with a paradox about the divergence between mathematical expectation and reasonable behavior.

\begin{quote}\itshape
The determination of the value of an item must not be based on its price, but rather on the utility it yields. The price of the item is dependent only on the thing itself and is equal for everyone; the utility, however, is dependent on the particular circumstances of the person making the estimate \citep{Bernoulli1738}.
\end{quote}

The St.\ Petersburg game offers a payoff of $2^n$ ducats if a fair coin first lands heads on flip $n$. The expected payoff is
\begin{equation}
\mathbb{E}[\text{payoff}] = \sum_{n=1}^{\infty} \frac{1}{2^n} \cdot 2^n = \sum_{n=1}^{\infty} 1 = \infty
\end{equation}
yet no reasonable person pays more than roughly 20 ducats to play. \citet{Bernoulli1738} resolved the paradox by proposing that agents maximize expected utility rather than expected monetary value. Under logarithmic utility $u(x) = \log(x)$, the expected utility for an agent with initial wealth $w$ is
\begin{equation}
\mathbb{E}[u] = \sum_{n=1}^{\infty} \frac{1}{2^n} \log\!\left(\frac{w + 2^n}{w}\right)
\end{equation}
which converges for any finite $w > 0$. The concavity of the logarithm ensures that the marginal utility of additional wealth diminishes, so the infinite sequence of doublings does not produce infinite expected utility.\footnote{Logarithmic utility resolves the original St.\ Petersburg paradox, but the paradox can be resurrected for any unbounded utility function by choosing a payoff sequence that grows fast enough. A \emph{super} St.\ Petersburg game paying $u^{-1}(2^n)$ ducats on flip $n$ restores infinite expected utility. Full resolution requires either bounded utility or additional decision-theoretic axioms.} The St.\ Petersburg game is the sharpest possible demonstration of the need for a reward function because its structure forces the issue: geometric probabilities paired with exponentially growing payoffs produce infinite expected monetary value, making any finite willingness to pay a direct refutation of expected value maximization. A game with merely large but finite expected payoff would show a quantitative gap between theory and behavior; the St.\ Petersburg game shows a qualitative one.

The distinction between raw payoff and its subjective transformation constitutes the first reward function in the modern sense. An agent optimizing monetary payoff and an agent optimizing a concave transformation of payoff will exhibit different behavior, and the choice between these objective functions is a modeling decision rather than a property of the environment. This is precisely the situation facing any reinforcement learning practitioner who must specify $r(s,a)$: the reward function encodes the designer's objectives, and different reward functions induce different optimal policies. Bernoulli's contribution was to recognize that the mapping from outcomes to values is not given by nature but chosen by the modeler. Chapter~\ref{section:rlhf} returns to this problem when reward functions are learned from human preferences rather than specified analytically.

\subsection{Bellman's Equation (1957)}

How should a warehouse manager decide today's order quantity when tomorrow's demand is uncertain and each day's stock depends on the previous day's decision? Multi-stage decision problems of this kind pervaded postwar operations research, and \citet{Bellman1957} recognized that they all share a common recursive structure.

A decision process consists of a state $s \in \mathcal{S}$, an action set $\mathcal{A}$, a transition kernel $P(s'|s,a)$, a reward function $r(s,a)$, and a discount factor $\gamma \in [0,1)$. The optimal value function satisfies the Bellman equation:
\begin{equation}
V^*(s) = \max_{a \in \mathcal{A}} \left[ r(s,a) + \gamma \sum_{s'} P(s'|s,a) \, V^*(s') \right]
\label{eq:bellman}
\end{equation}
The principle of optimality states that an optimal policy has the property that, whatever the initial state and initial decision, the remaining decisions must constitute an optimal policy with respect to the state resulting from the first decision. Value iteration computes $V^*$ by repeated application of the Bellman operator, converging at rate $\gamma$ per iteration.\footnote{Precisely, $\|V_{k+1} - V^*\|_\infty \leq \gamma \|V_k - V^*\|_\infty$. Achieving $\varepsilon$-optimal value estimates requires $k = \lceil \log(\|V_0 - V^*\|_\infty / \varepsilon) / \log(1/\gamma) \rceil$ iterations. For $\gamma$ near 1 the convergence is slow; with $\gamma = 0.99$ and initial error 100, reaching $\varepsilon = 0.01$ requires approximately 2{,}100 iterations.}

Bellman also identified the curse of dimensionality. If the state $s = (x_1, \ldots, x_d)$ with each component $x_i$ taking $m$ values, then $|\mathcal{S}| = m^d$, and both storage and computation grow exponentially in the number of state variables. This observation motivates every function approximation method in reinforcement learning.

\citet{BellmanDreyfus1962} applied dynamic programming to spacecraft trajectory optimization and missile allocation problems at the RAND Corporation, computing solutions on the JOHNNIAC computer. The trajectory problems involved continuous state spaces (position, velocity, fuel mass) discretized into grids, solving multi-stage optimization problems that the calculus of variations could not handle computationally. Bellman and Dreyfus claimed that dynamic programming ``resolved many problems seemingly inaccessible to the conventional methods of the calculus of variations,'' and the 1979 IEEE Medal of Honor citation for Bellman credited his invention of dynamic programming with opening ``a way for the application of sophisticated computer-oriented techniques in a wide variety of problem areas ranging from the design of guidance systems for space vehicles to pest control and network optimization.'' This was the largest-scale application of dynamic programming in the early 1960s, and the multi-dimensional continuous state spaces of the trajectory work foreshadowed the curse of dimensionality challenges that would later motivate function approximation.

In a three-page paper the same year, \citet{Bellman1958routing} applied the functional equation technique to a network routing problem: given $N$ cities connected by roads with varying travel times, find the minimum-time path between any two cities. The iterative algorithm converges in at most $N - 1$ steps, each of which relaxes every edge in the network, yielding a total complexity of $O(N^2)$ for dense graphs. The algorithm, now known as Bellman-Ford after independent work by Lester Ford at RAND, exploits the principle of optimality directly: the shortest path to any city must pass through a shortest sub-path to an intermediate city, so optimal paths can be built up from shorter optimal paths. A distributed variant of this algorithm became the original routing protocol for the ARPANET in 1969, the packet-switching network that evolved into the internet. The Routing Information Protocol (RIP), explicitly based on distributed Bellman-Ford, routed traffic across the early internet: each router maintained a distance vector to every destination and updated it by exchanging tables with its neighbors, precisely implementing Bellman's iterative relaxation in a decentralized setting. This is arguably the most widely deployed algorithm descended from dynamic programming, having routed every packet on the ARPANET and billions of packets daily during the internet's formative decades.

Around the same time, \citet{Scarf1960} used Bellman's recursive framework to prove a foundational result in inventory theory: under stochastic demand and a fixed ordering cost $K > 0$, the optimal inventory policy takes a simple $(s, S)$ threshold form. When inventory falls below $s$, order up to $S$; otherwise, do nothing. The proof established $K$-convexity of the value function by induction on the Bellman equation, a property preserved under the expectation and maximization operators.\footnote{A function $f$ is $K$-convex if $K + f(a+z) - f(a) - z[f(a) - f(a-b)]/b \geq 0$ for all $z > 0$, $b > 0$. $K$-convexity is preserved under (i) taking expectations with respect to random demand, and (ii) the dynamic programming operator that adds a fixed cost $K$ when ordering. The $(s,S)$ structure emerges because $K$-convex functions are minimized on intervals: it is optimal to do nothing when the current level is above a threshold.} Scarf's result was a theoretical proof rather than a computation, but it gave rigorous dynamic programming foundations to ordering rules that practitioners had long used heuristically, and it remains a cornerstone of the operations research curriculum. These early applications succeeded because their state spaces were naturally low-dimensional: the inventory problem has a single scalar state (stock level), and the trajectory problems had at most a handful of continuous variables (position, velocity, fuel mass) that could be discretized on a coarse grid. The recursive structure of the Bellman equation was transparent in these settings. Problems with high-dimensional or combinatorial state spaces, which would later motivate function approximation, were computationally out of reach.

\subsection{Howard's Policy Iteration (1960)}

Value iteration solves for the optimal value function directly, but it updates every state on every iteration and converges only asymptotically. \citet{howard1960} contributed an alternative algorithm, policy iteration, that alternates between two steps and converges in finitely many iterations. The evaluation step computes $V^\pi$ for the current policy $\pi$ by solving the linear system
\begin{equation}
V^\pi(s) = r(s, \pi(s)) + \gamma \sum_{s'} P(s'|s, \pi(s)) \, V^\pi(s') \quad \text{for all } s \in \mathcal{S}
\end{equation}
and the improvement step updates the policy via
\begin{equation}
\pi'(s) = \argmax_{a \in \mathcal{A}} \left[ r(s,a) + \gamma \sum_{s'} P(s'|s,a) \, V^\pi(s') \right]
\end{equation}
Howard proved that policy iteration converges to the optimal policy in a finite number of steps.\footnote{The number of deterministic policies is finite ($|\mathcal{A}|^{|\mathcal{S}|}$) and each iteration strictly improves the policy unless it is already optimal, so convergence is guaranteed. The worst-case number of iterations is exponential in $|\mathcal{S}|$, but empirically policy iteration converges in very few steps. For the automobile replacement problem, Howard reported convergence in three iterations. Ye (2011) proved a strongly polynomial bound of $O(|\mathcal{S}|^2 |\mathcal{A}| / (1 - \gamma))$ for the simplex method applied to the LP formulation.}

Howard devoted Chapter~5 of his monograph to three worked examples: taxicab routing across a small network of towns, a baseball batting order problem, and an automobile replacement problem. The automobile replacement example is the most instructive for economists: a fleet operator observes each vehicle's age (roughly 12 states corresponding to years of service), and must decide each period whether to keep the vehicle (paying maintenance costs that rise with age) or replace it (paying a fixed acquisition cost and resetting to new condition). Howard formulated this as a Markov decision process and showed that policy iteration found the cost-minimizing replacement rule in just a few iterations, cycling through only three policies before convergence. The automobile replacement problem was a natural showcase for policy iteration because its state space was small (roughly 12 age levels), the action was binary (keep or replace), and the Markov property held by construction: future costs depended only on the vehicle's current age, not on its maintenance history. These features made exact policy evaluation by solving a $12 \times 12$ linear system trivial, and the binary action set meant policy improvement reduced to a threshold comparison at each state. Howard reportedly developed the policy iteration algorithm while optimizing Sears, Roebuck \& Co.\ catalogue mailing decisions, choosing which customers to mail catalogues to based on purchase history. The automobile replacement example directly foreshadows the bus engine replacement model of \citet{Rust1987}, which we discuss in Chapter~\ref{section:rl_se}.

The Bellman equation and Howard's policy iteration together constitute the planning framework for sequential decisions under uncertainty. This framework assumes complete knowledge of $P$ and $r$. The challenge of learning optimal behavior without such knowledge is the subject of the remaining sections and the central problem of reinforcement learning. Chapter~\ref{section:planning_learning} develops the relationship between planning and learning in detail.

\subsection{Samuel's Checkers Program (1959)}

\citet{Samuel1959} set out to build a checkers program for the IBM 704 that could improve through experience rather than manual tuning, reasoning that learning from play should eventually reduce the burden of hand-coding evaluation logic. The program learned through self-play using two complementary mechanisms.

Rote learning stored every board position encountered along with its backed-up value, creating an expanding lookup table over the course of play. Generalization learning evaluated positions with a linear function
\begin{equation}
V(s) = \sum_{i} w_i \phi_i(s)
\end{equation}
where $\phi_i(s)$ were hand-crafted features: piece advantage, piece advancement, center control, mobility, and king-row threat. After each move, the weights were updated according to
\begin{equation}
\mathbf{w} \leftarrow \mathbf{w} + \alpha \left[ V(s_{t+1}) - V(s_t) \right] \nabla_{\mathbf{w}} V(s_t)
\end{equation}
This update rule adjusts the current position's value toward the value of the successor position, using the difference $V(s_{t+1}) - V(s_t)$ as a learning signal. The program reached expert-level play by 1962 and defeated Robert Nealey, a self-proclaimed expert, in a publicized match. Samuel's update is precisely TD(0) with linear function approximation, a connection formalized by \citet{sutton1988} twenty-nine years later. Checkers was well-suited to Samuel's approach because domain experts had already identified meaningful board features (piece advantage, center control, mobility), reducing the function approximation problem to learning a linear combination of informative inputs. The moderate branching factor ($b \approx 8$) kept the search tree manageable, and self-play generated unlimited training data at no cost. A game with less domain knowledge, a larger branching factor, or no natural features would have demanded the nonlinear function approximation that was still decades away.\footnote{Alpha-beta pruning, implemented following McCarthy's suggestion, reduces the minimax search complexity from $O(b^d)$ to $O(b^{d/2})$ under optimal move ordering, where $b$ is the branching factor and $d$ is the search depth. For checkers ($b \approx 8$, $d = 10$), this reduces the search space from $8^{10} \approx 10^9$ to $8^5 \approx 33{,}000$ positions.}

The linear polynomial assumed that features contributed independently to position value. In a follow-up study, Samuel replaced it with \emph{signature tables}: hierarchical lookup tables that captured within-group feature interactions. The 27 board parameters were divided into 9 groups of 3, and the evaluation became
\begin{equation}
V(s) = F\!\Bigl(T_1\bigl(\phi_1(s), \phi_2(s), \phi_3(s)\bigr),\; \ldots,\; T_9\bigl(\phi_{25}(s), \phi_{26}(s), \phi_{27}(s)\bigr)\Bigr)
\end{equation}
where each $T_g$ is a lookup table mapping a discretized feature triple to a score, and $F$ combines the group outputs through two further levels of table lookup. With parameters restricted to 3, 5, or 7 discrete values, the full three-level hierarchy required 20{,}956 table entries.\footnote{Each group of 3 parameters, discretized to $k \in \{3, 5, 7\}$ levels, yields a table of $k^3$ entries. With two further levels of table lookup combining group outputs, the total capacity was 20{,}956 entries stored in the IBM 7094's memory. This hierarchical decomposition can be viewed as a factored function approximation: it captures within-group interactions (up to 3-way) while assuming independence across groups.} This replacement of a linear combination with a hierarchical nonlinear mapping is a direct precursor to neural network function approximation: the conceptual step from Samuel's tables to Tesauro's hidden layers is replacing discrete lookups with continuous, differentiable activations.

Both the polynomial coefficients and the table entries were learned from a library of master-level games.\footnote{The library contained approximately 250{,}000 positions extracted from published master games. Each position was paired with the book-recommended move, providing a supervised training signal. By modern standards this is a small dataset; contemporary chess engines train on billions of self-play positions.} For each board parameter $i$, the program counted the number $H_i$ of non-book moves for which the parameter value exceeded that of the book-recommended move, and the number $L_i$ for which it fell below. The learned coefficient was
\begin{equation}
C_i = \frac{H_i - L_i}{H_i + L_i}
\end{equation}
a correlation measure that equals $+1$ when the parameter always correctly distinguishes the book move, $-1$ when it never does, and $0$ when it has no predictive power. This book learning procedure is supervised learning from expert demonstrations, not the self-play bootstrapping described above. Samuel used both: book learning to initialize coefficients, self-play to refine them. The distinction between learning from expert data and learning from self-generated experience anticipates the modern contrast between imitation learning and reinforcement learning.

\citet{Minsky1961} identified the credit assignment problem in his contemporary survey of artificial intelligence: in a game lasting many moves, which of the early decisions produced the eventual win or loss? Samuel's temporal difference update is one answer to this question, propagating information about outcomes backward through the sequence of positions without waiting for the game to end. The credit assignment problem remains central to reinforcement learning and is addressed explicitly by eligibility traces, policy gradients, and attention-based architectures.

\subsection{Sutton's TD($\lambda$) (1988)}

The gap between dynamic programming, which requires a complete model of the environment, and Monte Carlo estimation, which requires waiting for complete episodes, motivated a method that could learn incrementally from partial experience.

\begin{quote}\itshape
Whereas conventional prediction-learning methods are driven by the error between predicted and actual outcomes, TD methods are similarly driven by the error or difference between temporally successive predictions; with them, learning occurs whenever there is a change in prediction over time \citep{sutton1988}.
\end{quote}

\citet{sutton1988} formalized temporal difference learning and demonstrated its properties on a five-state random walk.

States A through E are arranged in a line with absorbing terminal states at each end. At each step the agent moves left or right with equal probability. The right terminal yields reward 1 and the left terminal yields reward 0. The true state values are $\frac{1}{6}, \frac{2}{6}, \frac{3}{6}, \frac{4}{6}, \frac{5}{6}$. Sutton compared TD(0), which updates after every transition, with Monte Carlo estimation, which updates only at the end of each episode. TD(0) converged faster and with less data than Monte Carlo on this problem. The five-state random walk is an ideal testbed for TD learning because it isolates the prediction problem in its purest form: there are no actions to choose, the transitions are purely stochastic, and the true state values are known analytically ($\frac{k}{6}$ for state $k$). This eliminates confounds from exploration, function approximation, and control, letting the comparison between TD and Monte Carlo focus entirely on how each method propagates value information.

The general TD($\lambda$) update combines both approaches through an eligibility trace:
\begin{equation}
V(s_t) \leftarrow V(s_t) + \alpha \, \delta_t \, e_t(s)
\end{equation}
where $\delta_t = r_{t+1} + \gamma V(s_{t+1}) - V(s_t)$ is the TD error and $e_t(s) = \gamma \lambda \, e_{t-1}(s) + \mathbbm{1}\{s = s_t\}$ is the eligibility trace for state $s$.\footnote{The eligibility trace formulation (backward view) is computationally efficient, requiring only $O(|\mathcal{S}|)$ storage per step. The equivalent forward view computes the $\lambda$-return $G_t^\lambda = (1-\lambda) \sum_{n=1}^{\infty} \lambda^{n-1} G_t^{(n)}$, a geometrically weighted average of $n$-step returns. The backward and forward views produce identical updates in the offline (batch) case; they differ slightly online due to the accumulating trace approximation.} Setting $\lambda = 0$ yields pure bootstrapping: the update uses only the immediate reward and the next state's estimated value. Setting $\lambda = 1$ recovers Monte Carlo returns, since the eligibility trace propagates the final outcome back through all visited states. Intermediate values of $\lambda$ interpolate between these extremes, and the optimal $\lambda$ is problem-dependent.

Updating from the next estimate rather than waiting for the final outcome allows learning to proceed at every timestep, even in continuing (non-episodic) tasks where no terminal state exists. This property makes TD methods applicable to a far broader class of problems than Monte Carlo.

\subsection{Q-Learning (1989)}

Temporal difference methods as developed by Sutton solve the prediction problem: estimating the value of a given policy. The control problem, finding the optimal policy without access to the transition dynamics, requires a different approach. \citet{WatkinsDayan1992}, formalizing Watkins's 1989 PhD thesis, provided an incremental method for dynamic programming that imposes limited computational demands:
\begin{equation}
Q(s_t, a_t) \leftarrow Q(s_t, a_t) + \alpha \left[ r_{t+1} + \gamma \max_{a'} Q(s_{t+1}, a') - Q(s_t, a_t) \right]
\end{equation}
The maximization over $a'$ is the critical feature. The update target uses the value of the best available action at the next state, regardless of which action the agent actually took to generate the data. This off-policy property separates the behavior policy (which may explore) from the target policy (which is greedy with respect to $Q$). The agent can follow an $\varepsilon$-greedy exploration strategy while learning about the optimal policy directly.\footnote{The $\max$ operator in the Q-learning target introduces a positive bias: $\mathbb{E}[\max_a Q(s,a)] \geq \max_a \mathbb{E}[Q(s,a)]$ when Q-values are noisy estimates. This maximization bias can cause systematic overestimation of action values. Double Q-learning (van Hasselt, 2010) addresses this by maintaining two independent Q-functions and using one to select and the other to evaluate, decoupling the max from the estimate.}

Convergence to $Q^*$ in the tabular setting requires two conditions: every state-action pair must be visited infinitely often, and the learning rate must satisfy the standard stochastic approximation conditions.\footnote{The conditions $\sum_t \alpha_t = \infty$ and $\sum_t \alpha_t^2 < \infty$ are the Robbins-Monro (1951) stochastic approximation conditions. The first ensures the step sizes are large enough in aggregate to reach the optimum from any initialization. The second ensures they shrink fast enough for the noise to average out. A common schedule satisfying both is $\alpha_t = c / (c + t)$ for constant $c > 0$. In practice, fixed small learning rates are often used instead, which do not guarantee convergence but track non-stationarity.} Under these conditions, $Q(s,a) \to Q^*(s,a)$ for all $(s,a)$ with probability one. The optimal policy is then $\pi^*(s) = \argmax_a Q^*(s,a)$. The convergence guarantee requires tabular representation and infinite visitation of every state-action pair, conditions naturally satisfied in small, discrete MDPs where an exploring agent revisits all cells. These conditions break down in large or continuous state spaces, where function approximation is necessary and the deadly triad (discussed below) becomes a concern.

The off-policy advantage is cleanly demonstrated by the cliff-walking problem \citep{sutton2018}. A $4 \times 12$ gridworld has a start cell in the bottom-left corner and a goal cell in the bottom-right corner, with a cliff running along the entire bottom row between them. Stepping onto the cliff yields a reward of $-100$ and returns the agent to start; all other steps yield $-1$. Q-learning, updating toward the greedy $\max_{a'} Q(s_{t+1}, a')$, learns the optimal shortest path directly along the cliff edge, because its target policy is greedy regardless of the exploratory actions actually taken. SARSA, which updates toward the action actually selected under the $\varepsilon$-greedy behavior policy, learns a longer safe path one row above the cliff, because its on-policy updates incorporate the cost of occasional exploratory steps off the edge. The two algorithms converge to different policies under the same $\varepsilon$-greedy exploration, isolating the consequence of off-policy versus on-policy learning.

Q-learning provided the first practical algorithm for model-free optimal control and became the foundation for a large family of value-based methods. Its extension to function approximation, however, would prove problematic, as discussed in the subsection on Baird's counterexample below.

\subsection{Williams' REINFORCE (1992)}

Value-based methods learn an action-value function and derive a policy from it. An alternative approach optimizes the policy directly by gradient ascent on expected return. \citet{williams1992} derived the policy gradient for a parameterized stochastic policy $\pi_\theta(a|s)$:
\begin{equation}
\nabla_\theta J(\theta) = \mathbb{E}_{\pi_\theta} \left[ \sum_{t=0}^{T} \nabla_\theta \log \pi_\theta(a_t|s_t) \, G_t \right]
\end{equation}
where $G_t = \sum_{k=0}^{T-t} \gamma^k r_{t+k+1}$ is the discounted return from time $t$. The log-derivative trick, $\nabla_\theta \pi_\theta = \pi_\theta \nabla_\theta \log \pi_\theta$, allows the gradient to be estimated from sampled trajectories without differentiating through the environment dynamics or the transition model. The gradient estimate is unbiased: in expectation over trajectories sampled from $\pi_\theta$, the sample gradient equals the true gradient $\nabla_\theta J(\theta)$.

The practical difficulty is variance. Because $G_t$ is a sum of stochastic rewards over the remainder of the episode, and because different trajectories may yield very different returns, the gradient estimates can be noisy. Variance reduction techniques, including baselines $b(s_t)$ subtracted from $G_t$ without introducing bias, are essential for practical performance.\footnote{The variance-minimizing baseline for the REINFORCE estimator is $b^*(s_t) = \mathbb{E}[G_t^2 \|\nabla_\theta \log \pi_\theta\|^2] / \mathbb{E}[\|\nabla_\theta \log \pi_\theta\|^2]$, which depends on the gradient magnitudes. In practice, $b(s_t) \approx V^\pi(s_t)$ is used, leading to the advantage actor-critic architecture where $G_t - b(s_t) \approx A^\pi(s_t, a_t)$.}

To see what this buys in practice, consider a robot arm that must apply a continuous torque $a \in \mathbb{R}$ to reach a target angle. Q-learning requires computing $\max_a Q(s,a)$ at every update, which is tractable when $a$ comes from a finite set but becomes a nested optimization problem when the action space is continuous. A policy gradient method sidesteps the issue entirely: parameterize the policy as a Gaussian $\pi_\theta(a|s) = \mathcal{N}(\mu_\theta(s),\, \sigma_\theta^2(s))$, sample an action, observe the return, and update $\theta$ by REINFORCE. The gradient $\nabla_\theta \log \pi_\theta(a|s)$ is available in closed form for the Gaussian, and no maximization over actions is required. This is why virtually all continuous-control results in reinforcement learning, from simulated locomotion to real robotic manipulation, descend from the policy gradient framework rather than from Q-learning.\footnote{Q-learning can be extended to continuous actions by training a separate optimization network or by restricting the Q-function to forms that admit closed-form maxima (e.g., NAF; Gu et al., 2016). These extensions are less general than policy gradients and introduce their own approximation errors.}

Policy gradients also handle problems where the optimal policy is genuinely stochastic. In a partially observed environment where distinct underlying states produce the same observation, a deterministic policy must take the same action in both states, which may be suboptimal. A stochastic policy can mix over actions and achieve higher expected return. \citet{SuttonMcAllester2000} emphasized this advantage, noting that the value-function approach is ``oriented toward finding deterministic policies, whereas the optimal policy is often stochastic.'' They also identified a subtler problem with value-based methods: an arbitrarily small change in estimated Q-values can cause the greedy policy to switch actions discontinuously, destabilizing learning with function approximation. Policy gradients avoid this because small changes in $\theta$ produce small changes in $\pi_\theta$.

\citet{SuttonMcAllester2000} demonstrated this concretely with a short corridor of four states where, in one interior state, the effect of the two actions is reversed: ``left'' moves right and ``right'' moves left. Any deterministic policy either always moves in the correct direction or always moves in the wrong direction at that state, and neither achieves the optimal expected return. A stochastic policy that mixes left and right with the correct probabilities does strictly better. An $\varepsilon$-greedy value-based method cannot discover this mixture because its policy is deterministic except for uniform random noise, whereas a policy gradient method parameterizes the action probabilities directly and converges to the optimal stochastic policy by gradient ascent on expected return.

\citet{SuttonMcAllester2000} proved the policy gradient theorem in the average reward setting, establishing that the gradient of expected return with respect to policy parameters takes the form $\nabla_\theta J(\theta) = \mathbb{E}_{s \sim d^\pi, a \sim \pi_\theta} [Q^\pi(s,a) \nabla_\theta \log \pi_\theta(a|s)]$, where $d^\pi$ is the stationary state distribution under $\pi$. The result is remarkable for what it excludes: the gradient does not require differentiating through $d^\pi$, even though $d^\pi$ depends on $\theta$ through the policy's effect on state transitions. Without this simplification, estimating the gradient from sampled trajectories would require knowing how the state distribution shifts with each parameter change, a quantity that is generally unavailable. The policy gradient theorem enabled the development of actor-critic methods that combine the low variance of value function estimation with the generality of policy parameterization.

\subsection{TD-Gammon (1994)}

The first large-scale demonstration that temporal difference learning with neural network function approximation could produce expert-level play was \citet{tesauro1994}'s TD-Gammon. Backgammon has approximately $10^{20}$ legal positions, far beyond the reach of tabular methods or exhaustive search. Backgammon was a uniquely favorable domain for this combination. The stochastic dice ensured that a greedy policy still explored diverse board positions, sidestepping the exploration problem that plagues deterministic games. The value function varied smoothly with board position (moving a single piece changes the win probability by a small amount), making it amenable to neural network approximation. And self-play generated unlimited on-policy training data, avoiding the off-policy component of the deadly triad.

TD-Gammon trained a feedforward neural network to estimate the probability of winning from any board position. The input $\mathbf{x}(s) \in \mathbb{R}^{198}$ encoded the board position using a unary representation.\footnote{The encoding used four binary units per board point per player: for each of the 24 points and each player, 4 binary units indicated occupancy of 1, 2, 3, or $\geq 4$ pieces. This yields $24 \times 2 \times 4 = 192$ units, plus 6 additional units for pieces on the bar and borne off, totaling 198. The unary encoding allows the network to learn nonlinear functions of piece count without requiring the hidden layer to first disentangle a scalar input.} A hidden layer of 80 sigmoid units (40 in version 2.0) fed a single sigmoid output:
\begin{equation}
\hat{V}(s) = \sigma\!\bigl(\mathbf{w}^\top \sigma(W\mathbf{x}(s) + \mathbf{b}) + c\bigr)
\end{equation}
where $\sigma$ is the logistic sigmoid, $W$ is the input-to-hidden weight matrix, and $\mathbf{w}$ is the hidden-to-output weight vector. The output $\hat{V}(s)$ estimates the probability of winning from position $s$.

The network was trained by self-play using TD($\lambda$) with $\lambda = 0.7$. After each move from position $s_t$ to $s_{t+1}$, the weights $\boldsymbol{\theta}$ (comprising $W$, $\mathbf{w}$, $\mathbf{b}$, and $c$) were updated by
\begin{equation}
\boldsymbol{\theta} \leftarrow \boldsymbol{\theta} + \alpha \bigl[\hat{V}(s_{t+1}) - \hat{V}(s_t)\bigr] \, \mathbf{e}_t
\end{equation}
where the eligibility trace $\mathbf{e}_t = \sum_{k=1}^{t} \lambda^{t-k} \nabla_{\boldsymbol{\theta}} \hat{V}(s_k)$ accumulates exponentially decayed gradients of past predictions, computed by backpropagation. At game's end, $\hat{V}(s_{t+1})$ is replaced by the outcome $z \in \{0, 1\}$, grounding the bootstrapping chain in actual wins and losses. At each turn, the program selected the move maximizing $\hat{V}(\mathrm{next}(s_t, a))$ over legal actions, a purely greedy policy. The stochastic dice rolls in backgammon provided natural exploration, a property Tesauro conjectured was critical to the method's success.\footnote{The stochastic dice in backgammon provide forced exploration: different rolls lead to different positions regardless of the player's strategy. This means a greedy policy still visits a diverse set of states. Deterministic games like chess and Go lack this property, which partly explains why TD-Gammon's pure self-play approach did not immediately transfer to those domains. AlphaGo Zero later solved the exploration problem in Go through MCTS-based action selection with added Dirichlet noise at the root node.}

TD-Gammon 1.0, trained on 300{,}000 self-play games with 1-ply lookahead, played at a strong intermediate level. Version 2.1, trained on 1.5 million games with 2-ply lookahead and 80 hidden units, achieved near-parity with Bill Robertie, a former world champion, losing by a single point over a 40-game session. The program discovered novel positional strategies that were subsequently adopted by the human backgammon community.

\begin{quote}\itshape
From an a priori point of view, this methodology appeared unlikely to produce any sensible learning, because random strategy is exceedingly bad, and because the games end up taking an incredibly long time \ldots\ The rather surprising result, after tens of thousands of training games, was that a significant amount of learning actually took place \citep{tesauro1994}.
\end{quote}

The success of TD-Gammon was surprising because it combined two individually unreliable components: bootstrapping (updating from estimates rather than outcomes) and function approximation (representing the value function with a neural network rather than a table). Researchers naturally asked whether these two components could be combined as freely in other settings. The answer, as the next subsection documents, was no.

\subsection{Baird's Counterexample and the Deadly Triad (1995)}

\citet{Baird1995} constructed a seven-state star MDP that demonstrated divergence of Q-learning with linear function approximation. The MDP has six ``outer'' states that all transition to a single ``inner'' state under the target policy. The off-policy behavior samples states uniformly. With linear function approximation, the weights grow without bound under repeated Q-learning updates, despite the simplicity of the MDP.

\begin{quote}\itshape
It is unfortunate that a reinforcement learning algorithm can be guaranteed to converge for lookup tables, yet be unstable for function-approximation systems that have even a small amount of generalization \citep{Baird1995}.
\end{quote}

The source of the instability is the interaction of three components: bootstrapping (updating from estimated values rather than observed returns), off-policy learning (training on data generated by a policy different from the target policy), and function approximation (representing the value function with a parameterized model rather than a table). \citet{sutton2018} later named this combination the deadly triad. Any two of the three components can be combined safely; it is the presence of all three simultaneously that permits divergence.

This result explains the asymmetry between Tesauro's success and Baird's failure. TD-Gammon used bootstrapping and function approximation but was on-policy: the training data came from the same self-play policy whose value was being estimated. Baird's counterexample used all three components, including off-policy updates, and diverged. Baird's star MDP was designed to isolate the failure mode with minimal structure: seven states, deterministic transitions, and a linear function approximation that cannot represent the true value function. The uniform off-policy sampling creates maximum mismatch with the target policy's visitation distribution, ensuring that the projection and Bellman operators interact adversarially. The example's simplicity is the point: divergence requires no exotic environment, only the simultaneous presence of the three ingredients. The deadly triad would remain the central theoretical obstacle in value-based deep reinforcement learning for two decades, until the stabilization techniques introduced by deep Q-networks partially addressed the problem in practice.\footnote{The divergence in Baird's example arises because the TD fixed point minimizes the projected Bellman error $\|V_\theta - \Pi T V_\theta\|$, where $\Pi$ projects onto the function approximation space. When the projection and the Bellman operator $T$ interact adversarially under off-policy sampling, the fixed point may not exist or may be unstable. Gradient TD methods (Sutton et al., 2009) resolve this by performing stochastic gradient descent on the mean-squared projected Bellman error, guaranteeing convergence at the cost of a second set of learned weights.}

\subsection{Deep Q-Networks (2015)}

End-to-end learning from raw sensory input became feasible when sufficient computation met architectural innovations for stabilizing the deadly triad. \citet{mnih2015} trained a single convolutional neural network to play 49 Atari 2600 games directly from pixel inputs ($210 \times 160 \times 3$) and a scalar score, using no game-specific features or hand-crafted representations. The network was trained by minimizing the squared temporal difference error:
\begin{equation}
L(\theta) = \mathbb{E}_{(s,a,r,s') \sim \mathcal{D}} \left[ \left( r + \gamma \max_{a'} Q(s', a'; \theta^-) - Q(s, a; \theta) \right)^2 \right]
\end{equation}
Two innovations stabilized learning. First, experience replay: transitions $(s, a, r, s')$ were stored in a replay buffer $\mathcal{D}$ and sampled uniformly for training, breaking the temporal correlation between consecutive updates that would otherwise destabilize gradient descent. Second, a target network: the bootstrap target used a frozen copy of the parameters $\theta^-$, updated to match $\theta$ only every $C$ steps.\footnote{The convolutional architecture used three layers (32 filters of $8 \times 8$ with stride 4, then 64 filters of $4 \times 4$ with stride 2, then 64 filters of $3 \times 3$ with stride 1) followed by a 512-unit fully connected layer. Input frames were downsampled to $84 \times 84$ grayscale and stacked 4 frames deep to provide velocity information. The replay buffer held $10^6$ transitions, the target network was updated every $C = 10{,}000$ steps, and training used RMSProp with a learning rate of $2.5 \times 10^{-4}$.} By holding the target fixed between updates, the regression target does not shift with each gradient step, reducing the feedback loop that drives divergence in the deadly triad.

DQN exceeded human-level performance on 29 of the 49 games using a single architecture and a single set of hyperparameters. The same network that learned to play Breakout (tracking a paddle to hit a ball) also learned to play Space Invaders (shooting descending aliens) and Montezuma's Revenge (navigating rooms with keys and doors, though poorly). The generality of the approach, learning from pixels with no domain knowledge, marked a qualitative shift in the field. The Atari 2600 platform was a natural proving ground for DQN because it provided a standardized interface (pixel frames in, one of 4--18 discrete actions out) across dozens of structurally distinct games, enabling a single architecture to be tested broadly. The discrete action spaces kept the $\max_{a'} Q(s',a')$ computation trivial, and the replay buffer was effective because Atari transitions, once conditioned on state and action, are approximately independent. Games requiring long-horizon planning or sparse rewards (such as Montezuma's Revenge) remained difficult, exposing the limits of value-based methods without hierarchical structure.

\subsection{AlphaGo and AlphaGo Zero (2016, 2017)}

The game of Go, with approximately $10^{170}$ legal positions and a branching factor of roughly 250, had long resisted both brute-force search and hand-crafted evaluation functions. \citet{Silver2016} combined learning and planning by integrating multiple components: a policy network trained by supervised learning on 30 million positions from human expert games, a second policy network improved by self-play reinforcement learning, a value network trained to predict game outcomes from positions, and Monte Carlo tree search (MCTS) guided by these networks. AlphaGo defeated Lee Sedol, one of the strongest Go players in history, four games to one in March 2016. Go was the right challenge for this architecture because its fixed $19 \times 19$ board is naturally suited to convolutional networks, its perfect information and deterministic transitions make MCTS's tree structure exact, and the game outcome provides a clean binary training signal. The enormous branching factor ($\sim$250) made pure search hopeless, ensuring that the neural network's role as a learned heuristic was essential rather than decorative.

\citet{Silver2017} removed human data entirely. AlphaGo Zero uses a single neural network $f_\theta(s) = (\mathbf{p}, v)$ that takes a board position $s$ as input and outputs a policy vector $\mathbf{p}$ over legal moves and a scalar value $v$ estimating the probability of winning. The training loop is self-contained: the program plays games against itself using MCTS guided by $f_\theta$, collects the MCTS-improved move probabilities $\boldsymbol{\pi}$ and the game outcome $z \in \{-1, +1\}$, then updates $f_\theta$ to minimize
\begin{equation}
\ell(\theta) = (z - v)^2 - \boldsymbol{\pi}^\top \log \mathbf{p} + c \|\theta\|^2
\end{equation}
where the first term is a value prediction loss, the second is a policy cross-entropy loss, and the third is $L_2$ regularization. After 72 hours of self-play on 4 TPUs,\footnote{AlphaGo Zero used a ResNet with 40 residual blocks (approximately 13 million parameters) and trained on 4.9 million self-play games. Each MCTS simulation performed 1{,}600 rollouts per move. The version that defeated Lee Sedol used a distributed system with 1{,}202 CPUs and 176 GPUs. The transition from AlphaGo to AlphaGo Zero demonstrated that removing human data and simplifying the architecture (single network instead of separate policy and value networks) improved both performance and training stability.} AlphaGo Zero surpassed all previous versions of AlphaGo, including the version that defeated Lee Sedol.

\citet{Igami2020} interpreted the AlphaGo architecture in econometric terms. The policy network is a conditional choice probability (CCP) estimator: given state $s$, it outputs a probability distribution over actions. The value network is a conditional value function (CVF) estimator: it maps states to expected future payoffs. In the language of structural econometrics, AlphaGo performs CCP estimation and forward simulation jointly, an observation that connects modern deep RL to the Hotz-Miller tradition in discrete choice. Chapter~\ref{section:planning_learning} develops this connection further.

\subsection{TRPO and PPO (2015, 2017)}

Policy gradient methods as described in the REINFORCE subsection suffer from a practical instability: a single large gradient step can move the policy into a region of parameter space where performance collapses, and recovery may be slow or impossible. The difficulty is that the objective $J(\theta)$ is not well approximated by its linear Taylor expansion when the step size is large relative to the curvature of the policy space. \citet{Schulman2015} addressed this problem with Trust Region Policy Optimization (TRPO), grounded in a monotonic improvement guarantee. For any two policies $\pi$ and $\tilde{\pi}$, the performance difference can be written exactly as
\begin{equation}
\eta(\tilde{\pi}) = \eta(\pi) + \mathbb{E}_{s \sim d^{\tilde{\pi}}, a \sim \tilde{\pi}} \left[ A^\pi(s,a) \right]
\end{equation}
where $\eta(\pi) = \mathbb{E}[\sum_t \gamma^t r_t]$ is the expected discounted return and $A^\pi(s,a) = Q^\pi(s,a) - V^\pi(s)$ is the advantage function. Because $d^{\tilde{\pi}}$ depends on the new policy and is difficult to compute, TRPO approximates it by using the surrogate objective
\begin{equation}
L_\pi(\tilde{\pi}) = \eta(\pi) + \mathbb{E}_{s \sim d^\pi, a \sim \tilde{\pi}} \left[ A^\pi(s,a) \right]
\end{equation}
which replaces $d^{\tilde{\pi}}$ with $d^\pi$. Schulman showed that the approximation error can be bounded:
\begin{equation}
\eta(\tilde{\pi}) \geq L_\pi(\tilde{\pi}) - C \, D_{\mathrm{KL}}^{\max}(\pi, \tilde{\pi})
\end{equation}
where $C$ is a constant depending on $\gamma$ and the advantage range,\footnote{The constant is $C = 4\gamma\epsilon / (1-\gamma)^2$ where $\epsilon = \max_{s,a} |A^\pi(s,a)|$, originally derived by Kakade and Langford (2002) for conservative policy iteration. The bound is loose: for $\gamma = 0.99$ and $\epsilon = 1$, $C \approx 400$, making the guaranteed improvement region extremely small. TRPO therefore ignores the bound and uses a fixed KL threshold $\delta \approx 0.01$, relying on empirical stability rather than the theoretical guarantee.} and $D_{\mathrm{KL}}^{\max}$ is the maximum KL divergence over states. Maximizing the right-hand side guarantees monotonic improvement. In practice, TRPO solves the constrained optimization problem
\begin{equation}
\max_\theta \; L_{\theta_{\text{old}}}(\theta) \quad \text{subject to} \quad \bar{D}_{\mathrm{KL}}(\theta_{\text{old}}, \theta) \leq \delta
\end{equation}
using the conjugate gradient method to compute the search direction and a line search to enforce the constraint. TRPO achieved strong results on continuous control tasks in MuJoCo (swimmer, hopper, walker) and on a subset of Atari games. Continuous control tasks in MuJoCo were a natural fit for trust region methods because the physics simulator produces smooth dynamics: small changes in torque produce small changes in trajectory, so small policy updates yield predictable performance changes. This smoothness is precisely the regime where a local surrogate objective closely approximates the true objective, making the trust region constraint informative rather than merely conservative.

A concrete example illustrates what the trust region buys. The MuJoCo Hopper is a one-legged robot with a 12-dimensional state (joint angles, velocities) and continuous torque controls. The reward is forward velocity minus a small effort penalty, with a bonus of $+1$ for staying upright; the episode terminates if the robot falls. \citet{Schulman2015} compared TRPO against the natural gradient method, which uses a fixed KL penalty coefficient instead of a constraint. On the Hopper, the natural gradient learned balanced standing but never walked forward, settling at a score of $-1$, the reward achievable without any forward motion. TRPO learned a hopping gait that made forward progress. The failure mode is instructive: a gradient step that is too large can shift the policy from a tentative walking gait into a regime where the hopper falls immediately. Once the policy falls, the collected data contains only falling trajectories, and the value estimates become unreliable, making recovery difficult. The KL constraint $\bar{D}_{\mathrm{KL}}(\theta_{\mathrm{old}}, \theta) \leq \delta$ prevents the policy from moving so far in one update that it enters this irrecoverable regime. The fixed penalty coefficient of the natural gradient cannot adapt to the local curvature of the policy space and either steps too conservatively (learning nothing) or too aggressively (falling).

\citet{Schulman2017} proposed Proximal Policy Optimization (PPO) as a simpler alternative that achieves comparable performance without the conjugate gradient computation. PPO replaces the KL constraint with a clipped surrogate objective:
\begin{equation}
L^{\text{CLIP}}(\theta) = \mathbb{E}_t \left[ \min\!\left( r_t(\theta) \hat{A}_t, \; \text{clip}(r_t(\theta), 1-\varepsilon, 1+\varepsilon) \hat{A}_t \right) \right]
\end{equation}
where $r_t(\theta) = \pi_\theta(a_t|s_t) / \pi_{\theta_{\text{old}}}(a_t|s_t)$ is the probability ratio and $\hat{A}_t$ is an estimator of the advantage function. The clipping removes the incentive for $r_t(\theta)$ to move outside the interval $[1-\varepsilon, 1+\varepsilon]$, effectively penalizing large policy updates without explicitly computing or constraining a divergence measure. PPO can be optimized by running multiple epochs of minibatch stochastic gradient descent on the same batch of trajectory data,\footnote{PPO typically runs 3--10 epochs of minibatch SGD per batch of trajectory data. The clipping threshold $\varepsilon = 0.2$ is standard. This multi-epoch reuse is unusual for on-policy methods (which traditionally discard data after one update) and is made safe by the clipping: once $r_t(\theta)$ exits $[1-\varepsilon, 1+\varepsilon]$, the gradient contribution from that sample is zeroed, preventing further policy drift on stale data.} making it substantially simpler to implement than TRPO.

In empirical comparisons, PPO outperformed A2C, TRPO, and the cross-entropy method on continuous control benchmarks and achieved the highest average reward on 30 of 49 Atari games among the methods tested. PPO became the default policy optimization algorithm for large-scale RL applications, including the preference learning systems discussed in the next subsection and in Chapter~\ref{section:rlhf}. The lesson from both TRPO and PPO is that constraining the magnitude of policy updates, whether through explicit KL constraints or surrogate clipping, is essential for stable optimization. Large policy changes are catastrophic because they move the agent into regions of state space where the current value estimates are unreliable.

\subsection{Reward Learning from Human Preferences (2017)}

All preceding methods assume that the reward function $r(s,a)$ is given. In many practical settings, however, the desired behavior is easy for a human to recognize but difficult to specify mathematically. \citet{Christiano2017} demonstrated that a reward function can be learned from a small number of pairwise human comparisons, bypassing the need for hand-coded reward specification.

\begin{quote}\itshape
Our algorithm fits a reward function to the human's preferences while simultaneously training a policy to optimize the current predicted reward function. We ask the human to compare short video clips of the agent's behavior, rather than to supply an absolute numerical score \citep{Christiano2017}.
\end{quote}

The method employs the Bradley-Terry preference model. Given two trajectory segments $\sigma^1 = (o_1^1, a_1^1, \ldots, o_k^1, a_k^1)$ and $\sigma^2 = (o_1^2, a_1^2, \ldots, o_k^2, a_k^2)$, the probability that a human evaluator prefers $\sigma^1$ to $\sigma^2$ is modeled as
\begin{equation}
P(\sigma^1 \succ \sigma^2) = \frac{\exp \sum_t \hat{r}_\phi(o_t^1, a_t^1)}{\exp \sum_t \hat{r}_\phi(o_t^1, a_t^1) + \exp \sum_t \hat{r}_\phi(o_t^2, a_t^2)}
\end{equation}
where $\hat{r}_\phi$ is a learned reward function parameterized by $\phi$. The reward model is trained by maximum likelihood on the comparison dataset, minimizing the cross-entropy loss between the predicted preference probabilities and the observed human labels. The policy is then optimized against $\hat{r}_\phi$ using PPO.

Approximately 900 pairwise comparisons sufficed to train a simulated MuJoCo humanoid to perform a backflip, a behavior that would be difficult to specify through a hand-coded reward function. The human evaluator watched pairs of short video clips and indicated which displayed more desirable behavior. No other reward signal was provided. The backflip task was well-chosen because the target behavior is visually distinctive and easy for a human to rank (more rotation is better) but difficult to specify as a mathematical reward function over joint angles and velocities. This asymmetry between recognizability and specifiability is the defining condition under which preference learning outperforms reward engineering. Tasks where the reward is easy to write (e.g., reaching a target position) gain little from the human-in-the-loop overhead.\footnote{The reward model $\hat{r}_\phi$ was trained as an ensemble of neural networks to provide uncertainty estimates. Each member was trained on a bootstrapped subsample of the comparison data. The predicted reward used for policy optimization was the ensemble mean, and disagreement among members served as a signal for which trajectory pairs to query next (active learning). The 900-comparison budget for the backflip task represents roughly 1 hour of human labeling time.}

The Bradley-Terry model is a conditional logit over trajectory pairs: the probability of choosing $\sigma^1$ over $\sigma^2$ is the softmax of the difference in summed rewards, identical in structure to the discrete choice models of \citet{McFadden1974}. The reward function $\hat{r}_\phi$ plays the role of the indirect utility function, and the pairwise comparison data are analogous to revealed preference observations. Reward hacking, in which the policy exploits features of $\hat{r}_\phi$ that diverge from the human evaluator's actual intent, remains a central open challenge. \citet{Ouyang2022} scaled the approach to natural language with InstructGPT, demonstrating that RLHF can align large language models with human preferences at scale. These issues are explored in Chapter~\ref{section:rlhf}.

\subsection{Toward a Unified Framework}

Every algorithm in this chapter performs some version of the same two operations: evaluate how good the current behavior is, then improve it. Bellman's equation and Howard's policy iteration made this loop explicit for the first time, alternating between computing $V^\pi$ and extracting a better $\pi'$ from it. \citet{sutton2018} call this loop Generalized Policy Iteration (GPI), and Chapter~\ref{section:planning_learning} develops the mathematics in full. What the fifteen developments above contribute, in retrospect, is a progressive relaxation of GPI's requirements at each stage, widening the class of problems the loop can solve.

Bernoulli supplied the objective: not raw payoffs but a subjective transformation of them, the reward function $r(s,a)$ that every subsequent method optimizes. Bellman supplied the recursive decomposition that makes evaluation tractable, and his routing paper showed the principle of optimality at work in the deterministic shortest-path case, an instance of GPI so clean that it became the backbone of internet routing. Howard made the two-step loop explicit: evaluate, then improve, then repeat. Thorndike and Rescorla-Wagner supplied the learning signal: behavior changes in proportion to the discrepancy between expected and received outcomes, the prediction error that Sutton would formalize as the TD error $\delta_t = r_{t+1} + \gamma V(s_{t+1}) - V(s_t)$. This error drives the evaluation step when the model is unknown: TD learning and Q-learning perform policy evaluation and improvement from samples rather than from $P$ and $r$, removing the requirement that the agent know its environment. Samuel and Tesauro showed that the evaluation step can use function approximation, replacing a table with a linear combination of features and then with a neural network, extending GPI to state spaces too large for exact computation. Baird showed that this extension is fragile: when bootstrapping, off-policy data, and function approximation interact, the evaluation step can diverge, breaking the loop entirely. Mnih's DQN stabilized it with experience replay and target networks, and Silver's AlphaGo combined a neural-network evaluation step with Monte Carlo tree search as the improvement step. REINFORCE and PPO reformulated the improvement step itself: instead of extracting a greedy policy from $Q$, they move the policy parameters directly by gradient ascent on expected return, with PPO constraining the step size to keep the loop stable. Christiano's preference learning removed the last fixed input: the reward function is no longer hand-specified but learned from human comparisons, so that even the objective of GPI is an output of the system.

The historical arc, viewed through GPI, is a progressive widening of what the agent is allowed to learn. Bellman and Howard assumed the model was known and computed everything by planning. TD methods and Q-learning learned the value function from experience but still required a given reward signal. Policy gradient methods learned the policy directly. Preference learning methods learned the reward function itself. At each stage, a quantity that was previously an input to GPI became an output of it. Recent theoretical work, surveyed in Chapter~\ref{section:planning_learning}, has made these connections mathematically precise: Q-learning is stochastic value iteration, natural policy gradient is soft policy iteration, and DQN with target networks satisfies the abstract conditions under which approximate dynamic programming converges. The common mathematical object is the Bellman operator and its contractive properties; the difference between planning and learning is whether that operator is applied with known transition probabilities or with sampled transitions. This same Bellman equation also underlies the structural econometric models of Chapter~\ref{section:rl_se}, where economists assume a model and solve for optimal behavior, while reinforcement learning agents learn that behavior from experience. The two traditions operate on the same recursive structure and increasingly use the same algorithms.
